\section{An enlarged monoidal model whose scalar defect disappears}
\label{sec:arch134-even-shift}

\subsection{ARCH-134-06: an exact even shift admits the tensor but normalizes every scalar}

\paragraph{Statement.}
Let $k$ be a field and $E$ an exact $k$-linear endofunctor of
$\operatorname{Vect}_k$, including the ultrapower quotient of
ARCH-134-05 when $k$ is finite.  Put
$\mathcal C=\operatorname{Ch}^{\ge0}(k)$, with its usual
\emph{cochain} tensor, and define the exact functor
\[
        (HC)^i=E(C^{i+2}),\qquad d_{HC}^i=E(d_C^{i+2})
                         \quad(i\ge0).
 \tag{ARCH134.S1}
\]
There is a symmetric monoidal square-zero monad
$\mathsf T=\mathrm{id}\oplus H$ on $\mathcal C$.
Its category $\mathcal A$ of algebras is complete and cocomplete,
abelian, $k$-linear, and has enough projectives.  It has a
symmetric monoidal $k$-bilinear tensor $\boxtimes$ which is right
exact in each variable.  An object is a cochain complex $C$ with
a chain map $a:HC\to C$ satisfying $aH(a)=0$.
The tensor unit and the scalar functor are
\[
       \mathbf1=k[0],\qquad \kappa(V)=(V[0],0)
                                  =\mathsf T(V[0]).
 \tag{ARCH134.S2}
\]
The scalar functor is exact, fully faithful and strong symmetric
monoidal, and its adjoints are
\[
 \varphi(C,a)=\operatorname{coker}(a^0:E(C^2)\to C^0),
 \qquad \psi(C,a)=\ker(d_C^0:C^0\to C^1).
 \tag{ARCH134.S3}
\]
For every vector space, with no finiteness restriction, the
actual derived scalar counit is an isomorphism:
\[
                L\varphi\kappa(V)=V[0].
 \tag{ARCH134.S4}
\]
Thus this construction repairs the tensor by changing the
two-sort model, but loses the intended scalar derived defect.

Moreover, $\mathcal M=\operatorname{Ch}_{\ge0}(\mathcal A)$,
now using an \emph{external homological} grading, satisfies the
scalar, normalization, tensor and interval predicates on
cote~134-8, physical pp.~48--49, footers~47--48.  For its actual
class $W$ of $L\varphi_{\mathcal M}$-equivalences one has
\[
                   \mathcal M[W^{-1}]\simeq D_{\ge0}(k).
 \tag{ARCH134.S5}
\]
This is therefore not a counterexample to the displayed
\emph{Thm(?)}.

\paragraph{Proof: the monad and its tensor.}
For any $k$-linear functor $E$, the maps
$E(U)\otimes V\to E(U\otimes V)$ are defined by
$x\otimes v\mapsto E(u\mapsto u\otimes v)(x)$.
They are bilinear, natural and associative with respect to
ordinary scalar tensoring.  Their symmetric counterparts are
defined in the same way.
Using the inclusion of each direct summand of a cochain tensor
gives natural chain maps
\[
 HC\otimes D\longrightarrow H(C\otimes D),\qquad
 C\otimes HD\longrightarrow H(C\otimes D).
 \tag{ARCH134.S6}
\]
For example, the first sends
$E(C^{i+2})\otimes D^j$ into
$E((C\otimes D)^{i+j+2})$.  The sign in the differential on
this summand is $(-1)^i=(-1)^{i+2}$, so it agrees with the
target differential.  The even shift is used here and also
in the symmetry identity.  Associativity of the maps follows
on each homogeneous component from associativity of the scalar
maps just defined.

Define the multiplication of $\mathsf T$ by
\[
 \mathsf T^2C=C\oplus HC\oplus HC\oplus H^2C
       \longrightarrow C\oplus HC,
       \qquad (c,x,y,z)\longmapsto(c,x+y),
 \tag{ARCH134.S7}
\]
and take the first inclusion as its unit.  This does not assume
$H^2=0$: it sets the multiplication of two $H$-operations to zero.
The monad identities follow by expanding $\mathsf T^3C$.
Terms containing no $H$ are unchanged; the three single-$H$
terms are added; every term containing at least two $H$'s is
killed by either composite.

On
$(C\oplus HC)\otimes(D\oplus HD)$ define the monoidal map
to $\mathsf T(C\otimes D)$ to be the ordinary inclusion on
$C\otimes D$, the two maps (ARCH134.S6) on the two single-$H$
summands, and zero on $HC\otimes HD$.  In the associativity
diagram, the no-$H$ component is ordinary associativity,
the single-$H$ components use the associativity already checked
for (ARCH134.S6), and all other components are zero on both
sides.  The unit identities use the normalized scalar maps.
Symmetry follows from the symmetry of (ARCH134.S6), including
the even-shift sign.  Compatibility with (ARCH134.S7) is checked
by the same expansion of
$\mathsf T^2C\otimes\mathsf T^2D$: each component with two
or more $H$'s is zero on both sides, and each single-$H$
component is the corresponding map (ARCH134.S6).  These
checks are all the monoidal-monad identities.

The underlying functor $\mathsf T$ is exact.  The base tensor
preserves cokernels in either variable, so
$\mathsf T(C\otimes-)$ and $\mathsf T(-\otimes D)$ preserve
reflexive coequalizers.  Seal, \emph{Tensors, monads and actions},
Theorem~2.6.4, therefore gives a symmetric monoidal tensor on
$\mathcal A$ with unit $\mathsf T(k[0])=k[0]$.
Its construction is the coequalizer of the two usual free
algebra maps
\[
 \mathsf T(\mathsf T C\otimes\mathsf T D)
       \rightrightarrows\mathsf T(C\otimes D)
       \longrightarrow (C,a)\boxtimes(D,b).
 \tag{ARCH134.S8}
\]
The two arrows use respectively the monad multiplication and
monoidal map, and the two algebra actions.  Proposition~2.5.2
and \S2.5.3 of that paper identify
\[
        \mathsf T C\boxtimes\mathsf T D
                         =\mathsf T(C\otimes D).
 \tag{ARCH134.S9}
\]
Only these stated existence and free-algebra results are used.

For completeness, bilinearity follows either from (ARCH134.S8)
or from its universal property for bilinear algebra maps:
all functors and arrows there are $k$-linear separately.
Right exactness can be checked from the same universal property.
A bilinear map out of $C\otimes D$ which kills
$C\otimes\operatorname{im}f$ factors uniquely through
$C\otimes\operatorname{coker}f$.  Its algebra compatibility
descends because $\mathsf T$ preserves the epimorphism induced
by $C\otimes D\to C\otimes\operatorname{coker}f$.
Thus the functor $C\boxtimes-$ preserves cokernels; the other
variable is identical.  This claim does not assert that
underlying vector-space coproducts are created by the forgetful
functor.

\paragraph{The abelian structure, projectives, and scalar adjoints.}
For an additive exact monad, kernels and cokernels of algebra
maps are created in the base category.  The induced action on
a kernel uses preservation of that kernel; the action on a
cokernel uses preservation of that cokernel.  Both algebra
identities follow by the corresponding monomorphism or
epimorphism.  Coimage equals image in the base, giving the
abelian assertion and exactness of the forgetful functor.
Limits are created by the forgetful functor, using their
universal properties.  Coproducts exist by writing every
algebra as its reflexive coequalizer of free algebras, taking
the coproducts of the free presentations, and then their
coequalizer.  Free algebras on any coproduct of base objects
exist, so this construction does not require $\mathsf T$ to
preserve coproducts.  Coproducts and cokernels give all colimits.

Let $D^n(V)$ denote the base cochain disk with $V$ in degrees
$n,n+1$ and differential $1$, for $n\ge0$.  It represents
$\operatorname{Hom}_k(V,C^n)$, so it is projective.
The base has enough projectives by taking sums of these disks.
The free algebra on a base projective is projective, since
the forgetful functor is exact.  A base projective epimorphism
$P\to UC$ gives an algebra epimorphism
$\mathsf T P\to\mathsf T UC\to C$; the last arrow is split
as a base map by the monad unit.  This proves enough
projectives in $\mathcal A$.

A chain map $(C,a)\to\kappa V$ has only its degree-zero
component, and the algebra condition says precisely that it
kills $a^0$.  In the other direction a chain map
$\kappa V\to(C,a)$ is a map $V\to\ker d_C^0$; its algebra
condition is automatic because $H(V[0])=0$.  This proves both
adjunctions (ARCH134.S3).  Exactness and full faithfulness of
$\kappa$ are immediate.  Formula (ARCH134.S9) proves its strong
symmetric monoidal compatibility with the scalar tensor.

\paragraph{The full scalar resolution.}
In the base category there is an exact augmented projective
resolution, in a separate homological resolution degree,
\[
 \cdots\longrightarrow D^2(V)\longrightarrow D^1(V)
       \longrightarrow D^0(V)\longrightarrow V[0]
       \longrightarrow0.
 \tag{ARCH134.S10}
\]
The map $D^n(V)\to D^{n-1}(V)$ is the identity in their common
cochain degree $n$ and zero elsewhere.  The augmentation is
the identity in cochain degree zero.  These are chain maps,
and exactness follows separately in each cochain degree.
Since $\mathsf T$ is exact and
$\mathsf T(V[0])=\kappa V$, applying it gives a projective
resolution in $\mathcal A$.
The action of the free algebra $\mathsf T C$ maps its copy
of $HC$ identically onto the second summand and kills $H^2C$.
Consequently
\[
     \varphi(\mathsf T C)=C^0,\qquad
     \varphi(\mathsf T D^n(V))=
       \begin{cases}V,&n=0,\\0,&n>0.\end{cases}
 \tag{ARCH134.S11}
\]
The augmented complex after $\varphi$ is therefore $V$ in
degree zero with its identity counit.  This proves (ARCH134.S4).
Although $\kappa V$ is free on the base object $V[0]$, that
base object is not projective when $V\ne0$: the epimorphism
$D^0(V)\to V[0]$ has no section.  We have used the explicit
resolution, not an assertion that the scalar objects are
projective.

\paragraph{The interval and actual localization.}
Give $\mathcal M=\operatorname{Ch}_{\ge0}(\mathcal A)$ its
external homological tensor formed from $\boxtimes$, finite
sums along each diagonal, and the external Koszul sign.
Its exact scalar functor and adjoints are
\[
 \varphi_{\mathcal M}=\varphi H_0,
 \qquad \kappa_{\mathcal M}V=\kappa V[0],
 \qquad \psi_{\mathcal M}=\psi\operatorname{ev}_0.
 \tag{ARCH134.S12}
\]
The checked finite-diagonal proof (ARCH134.10) applies to this
abelian category with enough projectives and identifies
$L\varphi_{\mathcal M}$ with $\mathbb L\varphi$ on external
complexes.  No exactness of arbitrary coproducts is used there.
Taking the external disk $J=D_1(\mathbf1)$ gives a natural
degreewise split monomorphism $X\to J\otimes X$.
The latter is the identity cone on $X$ and has its explicit
contracting homotopy.  Its $L\varphi_{\mathcal M}$ is therefore
zero.  Both tensoring with $J$ and tensoring with
$J/\mathbf1$ are exact, being respectively the identity cone
and a shift.  Formula (ARCH134.S4) gives the global unit
normalization with its actual counit.

Finally the exact right adjoint $\kappa$ gives the derived
adjunction
$\mathbb L\varphi:D_{\ge0}(\mathcal A)
\rightleftarrows D_{\ge0}(k):\kappa$.
Its counit is an isomorphism on every scalar complex, not only
on modules.  Indeed (ARCH134.S10) is functorial in $V$;
apply its free algebra resolution to each term of a connective
scalar complex and totalize the resulting first-quadrant
bicomplex.  Each total term is a finite sum of projectives.
The augmentation to the scalar complex is a quasi-isomorphism
by the exact resolution direction; after $\varphi$,
(ARCH134.S11) leaves literally the original scalar complex.
All diagonals used for these statements are finite.
An adjunction with invertible counit is a localization at the
maps inverted by its left adjoint, since the adjunction units
become isomorphisms and give inverse functors after localization.
Localizing the external chain category first at
quasi-isomorphisms gives (ARCH134.S5).  This proves every
assertion.\hfill$\square$

\paragraph{Reference and exact source use.}
Gavin J. Seal, \emph{Tensors, monads and actions},
Theory and Applications of Categories~28 (2013), no.~15,
403--434; Theorem~2.6.4, Proposition~2.5.2, \S2.5.3 and
Proposition~2.4.2.  The author's June~22, 2013 PDF has these
at pp.~21, 12--13 and 12 respectively.  The theorem and its
printed proof, the free-algebra proposition and its proof,
and the relevant definitions were inspected.  The more general
preceding tensor-coherence theorem is used through the
published result, not claimed re-proved here.
\url{https://www.epfl.ch/labs/hessbellwald-lab/wp-content/uploads/2018/10/MonadicActions.pdf}

\paragraph{Exact consequence and remaining gate.}
This gives a genuine enlarged monoidal construction with both
scalar adjoints and an interval.  Its scalar normalization
holds in every cardinality by (ARCH134.S10)--(ARCH134.S11),
so it removes the defect that the proposed counterexample
needed to retain.  Together with ARCH-134-05 it closes the
specified quotient/strength and exact-shift attempts.  It
does not decide whether some different enlarged tensor on
the original two-sort category exists, whether that category
can give a full-sheet counterexample, or whether the bare
\emph{Thm(?)} always holds.  The arbitrary-ring raw ordinary
$G$ counit SM25.2 is also unchanged.

\paragraph{Verification status.}
The author completed a whole saved-proof reading: the even-shift
chain signs, every monoidal-monad summand, the exact published
tensor hypotheses, the abelian and coproduct constructions, both
scalar adjunctions, the full projective disk resolution, and the
finite-diagonal interval/localization argument were checked.  No
gap was found.  Reviewer cote123 subsequently completed a full
saved-proof reading, re-read (ARCH134.10), inspected the stated
Seal inputs directly, and reported PASS.  Root completed a whole
saved-proof reading and directly checked Seal's Theorem~2.6.4,
Proposition~2.5.2, their printed proofs and definitions, and
reported PASS.  The actual root event is in the central register;
it is not counted as an independent re-proof of the inherited
derived-comparison result.
